\documentclass[11pt,a4paper]{article}
\usepackage[margin=2.2cm]{geometry}
\usepackage[T1]{fontenc}
\usepackage{lmodern}
\usepackage{xcolor}
\usepackage{booktabs}
\usepackage{longtable}
\usepackage{enumitem}
\usepackage{hyperref}
\usepackage{microtype}
\hypersetup{colorlinks=true,linkcolor=blue!55!black,urlcolor=blue!65!black}
\setlist{nosep,leftmargin=*}

\title{Pilot Fabric 0.32\\Trusted Private-Phone HTTPS}
\author{Tessaris}
\date{30 August 2026}

\begin{document}
\maketitle

\section{Purpose}

Browser camera APIs require a secure context.  Pilot Fabric 0.32 therefore adds
a TLS private companion without turning a public cloud into the home's trust
anchor.  The mother brain creates a household certificate authority (CA), issues
a LAN server certificate and runs the existing tokenized controller over HTTPS.
The original HTTP page remains a local bootstrap and recovery path.

\section{Local public-key infrastructure}

Certificate generation uses the cryptography library already included in the
Pilot package.  It does not execute a shell command, contact a CA service, or
upload a key.  The household root uses a 3072-bit RSA key and SHA-256 signature,
is constrained to CA use with path length zero, and has key-signing and revocation-
list usage only.  The server uses a separate 3072-bit RSA key and a 397-day
certificate restricted to TLS server authentication.

The leaf Subject Alternative Name contains the active private LAN address, the
mother's local hostname, localhost and the loopback address.  If DHCP changes the
LAN address, Pilot issues a new leaf from the same household root.  This keeps the
trusted root stable while ensuring the certificate matches the actual connection
target.

\section{Key custody}

The TLS directory and private keys use owner-only filesystem permissions.  The
CA key and server key never enter an API response, QR code, television projection,
status payload, package, or audit record.  Only the public CA certificate may be
downloaded.  Public status contains certificate fingerprints, current address,
hostname and the explicit-trust requirement.

\section{Two-port onboarding model}

\begin{longtable}{p{0.23\linewidth}p{0.20\linewidth}p{0.49\linewidth}}
\toprule
\textbf{Port} & \textbf{Role} & \textbf{Behavior} \\
\midrule
\endhead
8767 & Bootstrap and recovery & Shows the rotating pairing code, public CA
download, secure-controller link, fingerprint and trust instructions. \\
8769 & Trusted controller & Serves the same private controller and perception
endpoints using TLS 1.2 or later after the phone trusts the household root. \\
\bottomrule
\end{longtable}

A correct pairing code redirects to the secure controller when TLS is available.
Both routes retain the secret path token, CSRF token, request-size limits and IP
rate limits.  HTTPS adds HSTS.  Responses remain non-cacheable and deny framing,
cross-origin resources and microphone access; camera permission is restricted to
the controller's own origin.

\section{Explicit Apple trust flow}

Pilot displays the root SHA-256 fingerprint on the trusted laptop dashboard and
on the pairing page.  The owner downloads the public certificate, reviews the
profile, compares the fingerprint, installs it and then explicitly enables full
trust in Certificate Trust Settings.  These are platform security decisions and
Pilot never clicks, suppresses or automates them.

Apple documents that manually installed roots on unsupervised iPhone and iPad
devices are not automatically trusted for SSL/TLS.  Full trust must be enabled
separately.  Apple recommends Configurator or mobile-device management for managed
organizational deployment.  Consumer Pilot preserves manual control and allows
trust to be revoked by removing the installed profile.

\section{Camera and observer behavior}

Once the certificate is trusted, the HTTPS page is a secure browser context and
may call \texttt{get\allowbreak User\allowbreak Media} after the user grants camera permission.  Observer
Mode remains visible, ten-minute limited, persona-bound and frame-rate limited.
Page hiding stops tracks and timers; return requires \emph{Resume after phone
sleep}.  TLS enables the API but does not weaken consent or permit invisible
background capture.

\section{Failure and recovery}

If CA creation, leaf issuance or TLS binding fails, Fabric continues to expose
the HTTP recovery controller and records only the failure class.  No private key
material appears in logs.  The HTTP controller remains suitable for remote and
approval operations, while the UI honestly explains that continuous camera
access needs HTTPS or a signed native client.

\section{Verification}

Automated tests verify owner-only key modes, required IP and DNS SAN values, root
reuse, leaf rotation after an address change, TLS chain validation against the
generated root, HSTS, public-certificate delivery and the absence of private-key
delivery.  The complete Fabric suite rechecks CSRF, rate limits, controller
commands, observer leases, local frame deletion, multi-frame stability, verified
navigation, webOS operation and audit integrity.

\section{Remaining work}

A native signed iOS and Android companion would remove the manual-root onboarding
step, use OS key stores and support clearer background lifecycle behavior.  The
next data-plane stage is authenticated provider metadata for programme, episode,
entitlement and playback position where each service permits it.  Such adapters
must remain persona-bound and must never export session cookies to TV nodes.

\section{References}

\begin{itemize}
  \item Apple, \href{https://support.apple.com/en-us/102390}{Trust manually
  installed certificate profiles in iOS, iPadOS, and visionOS}.
  \item Mozilla, \href{https://developer.mozilla.org/en-US/docs/Web/API/MediaDevices/getUserMedia}{MediaDevices.getUserMedia secure-context requirements}.
\end{itemize}

\end{document}
